% This must be in the first 5 lines to tell arXiv to use pdfLaTeX, which is strongly recommended.
\pdfoutput=1
% In particular, the hyperref package requires pdfLaTeX in order to break URLs across lines.

\documentclass[11pt]{article}

% Remove the "review" option to generate the final version.
% \usepackage[review]{acl}
\usepackage[]{acl}

% Standard package includes
\usepackage{times}
\usepackage{latexsym}

% For proper rendering and hyphenation of words containing Latin characters (including in bib files)
\usepackage[T1]{fontenc}
% For Vietnamese characters
% \usepackage[T5]{fontenc}
% See https://www.latex-project.org/help/documentation/encguide.pdf for other character sets

% This assumes your files are encoded as UTF8
\usepackage[utf8]{inputenc}

% This is not strictly necessary, and may be commented out,
% but it will improve the layout of the manuscript,
% and will typically save some space.
\usepackage{microtype}

% This is also not strictly necessary, and may be commented out.
% However, it will improve the aesthetics of text in
% the typewriter font.
\usepackage{inconsolata}

\usepackage{authblk}
\usepackage{algorithm}
\usepackage{algorithmic}

\usepackage{amsmath}
\usepackage{times}
\usepackage{latexsym}
\usepackage{booktabs}
\usepackage{tabularx}
\usepackage{tikz}
\usepackage{pgfplots}
\usepackage{pgfplotstable}
\usepackage{enumitem}

% Standard package includes
\usepackage{soul}
\usepackage{pifont}
\usepackage{xcolor,colortbl}

\usepackage{cleveref}
\usepackage[]{mdframed}
\usepackage{framed}
\newcommand{\cmark}{\ding{51}}%
\newcommand{\xmark}{\ding{55}}%

\usepackage{url}
\usepackage{graphicx}

\usepackage{tabularx}
\usepackage{xcolor}
\usepackage{makecell}
\usepackage{enumitem}

\definecolor{Gray}{gray}{0.9}

% If the title and author information does not fit in the area allocated, uncomment the following
%
%\setlength\titlebox{<dim>}
%
% and set <dim> to something 5cm or larger.

\title{To Tell The Truth:\\Language of Deception and Language Models}
% \author { 
%     Bodhisattwa Prasad Majumder\textsuperscript{\rm 1}\thanks{Authors contributed equally},
%     Sanchaita Hazra\textsuperscript{\rm 2}\textsuperscript{\rm *}
% }
% \affiliations {
%     \textsuperscript{\rm 1} Allen Institute for AI\\
%     \textsuperscript{\rm 2} University of Utah\\
%     \texttt{bodhisattwam@allenai.org, sanchaita.hazra@utah.edu}
% }


\author[$1$]{\textbf{Sanchaita Hazra}*}
\author[$2$]{\textbf{Bodhisattwa Prasad Majumder}*}
\affil[$1$]{University of Utah}
\affil[$2$]{Allen Institute for AI \protect \\ \texttt{sanchaita.hazra@utah.edu, bodhisattwam@allenai.org} \protect \\
\small \texttt{*contributes equally}}

\setlength{\affilsep}{0em}


\begin{document}
\maketitle
\begin{abstract}
Text-based false information permeates online discourses, yet evidence of people's ability to discern truth from such deceptive textual content is scarce. We analyze a novel TV game show data where conversations in a high-stake environment between individuals with conflicting objectives result in lies. We investigate the manifestation of potentially verifiable language cues of deception in the presence of objective truth, a distinguishing feature absent in previous text-based deception datasets. We show that there exists a class of detectors (algorithms) that have similar truth detection performance compared to human subjects, even when the former accesses only the language cues while the latter engages in conversations with complete access to all potential sources of cues (language and audio-visual). Our model, built on a large language model, employs a bottleneck framework to learn discernible cues to determine truth, an act of reasoning in which human subjects often perform poorly, even with incentives. Our model detects novel but accurate language cues in many cases where humans failed to detect deception, opening up the possibility of humans collaborating with algorithms and ameliorating their ability to detect the truth.
\end{abstract}

\input{files/introduction}
\input{files/related_work}
\input{files/data}
\input{files/method}
\input{files/experiments}
\input{files/learning}
\input{files/conclusion}


% Entries for the entire Anthology, followed by custom entries
\bibliography{anthology,custom}

\appendix
\input{files/appendix}

% \section{Example Appendix}
% \label{sec:appendix}

% This is an appendix.

\end{document}
